% Visual analysis at baseline and follow-up: what was read. The cohort criteria and outcome are in
% methodology/02_population_outcome.tex (sec:method-population); where incident plaque developed is in
% results/00_cohort.tex (sec:results-cohort), moved there 2026-10-31.
\section{Visual Analysis and Patient Characteristics}
\label{sec:data-plaque-location}

Every baseline and follow-up CCTA was analysed visually.
\clarify{how and by whom the visual analysis was made.}
For each of the 17 coronary segments, the reader recorded the stenosis grade (none, minimal, mild, moderate, severe, or occlusion), the plaque type (calcified, mostly calcified, mostly non-calcified, or non-calcified), high-risk plaque features, and evaluability.
The reads were summarised per participant in the segment involvement score (SIS) and the segment stenosis score (SSS), counting the segments with plaque and summing the stenosis grades, respectively.
In the delivered baseline table, segments without a grade are filled with none, and whether a scan was read is recorded in a separate flag; all 1,801 baseline scans were read.

All SCORE2 risk factors \cite{score2_2021} were also recorded at both visits: age, sex, current smoking, systolic blood pressure, and total and high-density lipoprotein (HDL) cholesterol, from which non-HDL cholesterol is derived.
